% Chapter 2: Smooth Maps
% Lee, Introduction to Smooth Manifolds (2nd ed., GTM 218).
% Generated from the section extraction; nodes carry only Lean that is
% verified sorry-free. See ../../UPSTREAM_LEAN_AUDIT.md.


\section{Smooth Functions and Smooth Maps}

\begin{remark}
\label{rem:2-8-extra-1}
\lean{ContMDiffMap}

\mathlibok
The term \emph{function} denotes a real- or vector-valued map, while \emph{map} may have arbitrary codomain.
\end{remark}

\begin{notation}
\label{not:2-8-extra-3}
\lean{C}

\mathlibok
For a smooth manifold $M$, write $C^\infty(M)$ for the real-valued smooth
functions on $M$; with pointwise operations it is a commutative associative
$\mathbb R$-algebra.
\end{notation}


\begin{lemma}[Smooth Coordinate Representations]
\label{exer:2-3}
\lean{smooth_chart_representation_of_contMDiff}
\leanok


\dcref{ch2:2.3}
If $f:M\to\mathbb R^k$ is smooth and $(U,\varphi)$ is any smooth chart, then $f\circ\varphi^{-1}:\varphi(U)\to\mathbb R^k$ is smooth.
\end{lemma}

\begin{definition}
\label{def:2-8-extra-4}
\lean{writtenInExtChartAt}

\mathlibok
For $f:M\to\mathbb R^k$ and a chart $(U,\varphi)$, its coordinate representation is $\widehat f=f\circ\varphi^{-1}:\varphi(U)\to\mathbb R^k$. The function $f$ is smooth exactly when these representations are smooth in a chart around every point; by Lemma~\ref{exer:2-3}, smoothness then holds in every smooth chart.
\end{definition}

\begin{definition}
\label{def:2-8-extra-5}
\lean{ContMDiff}

\mathlibok
For smooth manifolds $M,N$, a map $F:M\to N$ is smooth if for every $p\in M$ there are smooth charts $(U,\varphi)$ about $p$ and $(V,\psi)$ about $F(p)$ with $F(U)\subseteq V$ such that $\psi\circ F\circ\varphi^{-1}:\varphi(U)\to\psi(V)$ is smooth. For manifolds with boundary, maps from subsets of $\mathbb H^n$ are smooth when locally extendible to smooth maps, and maps into subsets of $\mathbb H^n$ are smooth as maps into $\mathbb R^n$. The function definition is the case $N=\mathbb R^k$ with $\psi=\operatorname{Id}$.
\end{definition}

\begin{proposition}
\label{prop:2-4}
\lean{ContMDiff.continuous}

\mathlibok
\dcref{ch2:2.4}
Every smooth map is continuous.
\end{proposition}

\begin{proof}
\sketch For a smooth $F:M\to N$ and each $p$, choose charts as in the definition. The local expression
\[
F|_U=\psi^{-1}\circ(\psi\circ F\circ\varphi^{-1})\circ\varphi
\]
is continuous, hence $F$ is continuous on $M$.
\end{proof}

\begin{proposition}[Equivalent Characterizations of Smoothness]
\label{prop:2-5}
\lean{contMDiff_iff}

\mathlibok
\dcref{ch2:2.5}
For a map $F:M\to N$ between smooth manifolds (with or without boundary), $F$ is smooth if and only if either:

(a) For every $p\in M$ there are smooth charts $(U,\varphi)$ about $p$ and $(V,\psi)$ about $F(p)$ such that $U\cap F^{-1}(V)$ is open and $\psi\circ F\circ\varphi^{-1}$ is smooth on $\varphi(U\cap F^{-1}(V))$; or

(b) $F$ is continuous and there are smooth atlases $\{(U_\alpha,\varphi_\alpha)\}$ and $\{(V_\beta,\psi_\beta)\}$ for $M,N$ such that every $\psi_\beta\circ F\circ\varphi_\alpha^{-1}$ is smooth on $\varphi_\alpha(U_\alpha\cap F^{-1}(V_\beta))$.
\end{proposition}

\begin{proposition}[Smoothness Is Local]
\label{prop:2-6}
\lean{contMDiff_of_locally_contMDiffOn}

\mathlibok
\dcref{ch2:2.6}
For a map $F:M\to N$ between smooth manifolds, if every $p\in M$ has a neighborhood on which $F$ is smooth, then $F$ is smooth. Conversely, the restriction of a smooth map to every open subset is smooth.
\end{proposition}


\begin{corollary}[Gluing Lemma for Smooth Maps]
\label{cor:2-8}
\lean{gluing_lemma_for_smooth_maps}
\uses{prop:2-6}



\dcref{ch2:2.8}
If $\{U_\alpha\}_{\alpha\in A}$ is an open cover of $M$ and smooth maps $F_\alpha:U_\alpha\to N$ agree on every overlap, there is a unique smooth $F:M\to N$ with $F|_{U_\alpha}=F_\alpha$.
\end{corollary}

\begin{definition}
\label{def:2-8-extra-6}
\lean{extended_coordinate_representation_mapsTo}
\leanok


For a smooth $F:M\to N$ and smooth charts $(U,\varphi)$, $(V,\psi)$, define $\widehat F=\psi\circ F\circ\varphi^{-1}$. It maps $\varphi(U\cap F^{-1}(V))$ into $\psi(V)$.
\end{definition}

\begin{lemma}[Chart Independence of Smooth Coordinate Representations]
\label{exer:2-9}
\lean{contDiffOn_coord_repr_of_contMDiff}
\leanok


\dcref{ch2:2.9}
The coordinate representation of a smooth map with respect to every pair of smooth charts is smooth.
\end{lemma}

\begin{proposition}
\label{prop:2-10}
\lean{contMDiff_const}
\uses{exer:2-9}

\mathlibok

\dcref{ch2:2.10}
For smooth manifolds $M,N,P$: every constant map $M\to N$, the identity of $M$, and the inclusion of an open submanifold $U\hookrightarrow M$ are smooth; if $F:M\to N$ and $G:N\to P$ are smooth, then $G\circ F$ is smooth.
\end{proposition}

\begin{proof}
\sketch Parts (a)--(c) have constant, identity, or coordinate-inclusion representations. For (d), choose a chart $(V,\theta)$ around $F(p)$ on which $G$ has image in a chart $(W,\psi)$, then choose $(U,\varphi)$ with $U\subseteq F^{-1}(V)$. The identity
\[
\psi\circ(G\circ F)\circ\varphi^{-1}=(\psi\circ G\circ\theta^{-1})\circ(\theta\circ F\circ\varphi^{-1})
\]
is a composition of smooth Euclidean maps.
\end{proof}


\begin{proposition}
\label{prop:2-12}
\lean{contMDiff_pi_iff}
\uses{prob:2-2}
\leanok


\dcref{ch2:2.12}
Let $M_1,\ldots,M_k,N$ be smooth manifolds with or without boundary, with at most one $M_i$ having nonempty boundary. For projections $\pi_i:M_1\times\cdots\times M_k\to M_i$, a map $F:N\to M_1\times\cdots\times M_k$ is smooth if and only if each $F_i=\pi_i\circ F$ is smooth.
\end{proposition}

\begin{proof}
\leanok
\sketch In product charts, the coordinate representation of $F$ is the tuple
of the coordinate representations of the $F_i$. Such a Euclidean map is smooth
if and only if each component is smooth.
\end{proof}

\begin{remark}
\label{rem:2-8-extra-7}
\lean{OpenPartialHomeomorph}

\mathlibok
Unless specified otherwise, a function or map is set-theoretic, a manifold is topological, and a coordinate chart is a homeomorphism to an open subset of $\mathbb R^n$ or $\mathbb H^n$. Continuity or smoothness is imposed only when stated; saying that $F:M\to N$ is smooth implicitly gives smooth structures on $M$ and $N$.
\end{remark}

\begin{example}
\label{ex:2-13}
\lean{manifold_pi_projection_contMDiff}
\uses{ex:1-33, ex:1-8, prob:1-8, prop:2-12}
\leanok



\dcref{ch2:2.13}
(a) Every map from a zero-dimensional manifold is smooth.

(b) For the standard circle, $\varepsilon:\mathbb R\to\mathbb S^1$, $\varepsilon(t)=e^{2\pi it}$, has angle-coordinate representation $2\pi t+c$ and is smooth.

(c) The map $\varepsilon^n:\mathbb R^n\to\mathbb T^n$, $\varepsilon^n(x^1,\ldots,x^n)=(e^{2\pi ix^1},\ldots,e^{2\pi ix^n})$, is smooth by Proposition~\ref{prop:2-12}.

(d) The inclusion $\iota:\mathbb S^n\hookrightarrow\mathbb R^{n+1}$ has graph-coordinate representation
\[
(u^1,\ldots,u^n)\longmapsto(u^1,\ldots,u^{i-1},\pm\sqrt{1-|u|^2},u^i,\ldots,u^n),
\]
which is smooth for $|u|^2<1$.

(e) The quotient map $\pi:\mathbb R^{n+1}\setminus\{0\}\to\mathbb{RP}^n$ has, on the $i$th projective chart, representation
\[
(x^1,\ldots,x^{n+1})\longmapsto(x^1/x^i,\ldots,x^{i-1}/x^i,x^{i+1}/x^i,\ldots,x^{n+1}/x^i),
\]
so it is smooth.

(f) The restriction $q=\pi\circ\iota:\mathbb S^n\to\mathbb{RP}^n$ is smooth.

(g) Each product projection $\pi_i:M_1\times\cdots\times M_k\to M_i$ is smooth.
\end{example}

\section{Diffeomorphisms}

\begin{definition}
\label{def:2-9-extra-1}
\lean{M}

\mathlibok
For smooth manifolds $M,N$, a diffeomorphism is a smooth bijection $F:M\to N$ whose inverse is smooth. The manifolds are diffeomorphic, written $M\approx N$, when such a map exists.
\end{definition}

\begin{example}
\label{ex:2-14}
\lean{smoothChartDiffeomorph}
\leanok


\dcref{ch2:2.14}
(a) The maps
\[
F(x)=\frac{x}{\sqrt{1-|x|^2}}:\mathbb B^n\to\mathbb R^n,
\qquad G(y)=\frac{y}{\sqrt{1+|y|^2}}:\mathbb R^n\to\mathbb B^n
\tag{2.1}
\]
are smooth inverses, so $\mathbb B^n\approx\mathbb R^n$.

(b) Every smooth chart $\varphi:U\to\varphi(U)\subseteq\mathbb R^n$ is a diffeomorphism.
\end{example}

\begin{theorem}[Diffeomorphism Invariance of Dimension]
\label{thm:2-17}
\lean{diffeomorphic_dimension_eq}
\leanok


\dcref{ch2:2.17}
Nonempty smooth manifolds of dimensions $m$ and $n$ can be diffeomorphic only when $m=n$.
\end{theorem}

\begin{proof}
\leanok
\sketch A diffeomorphism restricts in local charts to a diffeomorphism between open subsets of $\mathbb R^m$ and $\mathbb R^n$. Invariance of Euclidean dimension (Proposition C.4) gives $m=n$.
\end{proof}

\section{Partitions of Unity}

\begin{lemma}
\label{lem:2-20}
\lean{expNegInvGlue.contDiff}

\mathlibok
\dcref{ch2:2.20}
The function
\[
f(t)=\begin{cases}e^{-1/t},&t>0,\\0,&t\leq0\end{cases}
\]
is smooth on $\mathbb R$.
\end{lemma}

\begin{proof}
\sketch For every integer $k\geq0$, $\lim_{t\searrow0}e^{-1/t}/t^k=0$. Inductively, for $t>0$,
\[
f^{(k)}(t)=p_k(t)\frac{e^{-1/t}}{t^{2k}},
\]
where $p_k$ is a polynomial of degree at most $k$; differentiating gives the same form with $p_{k+1}(t)=t^2p_k'(t)+p_k(t)-2ktp_k(t)$. The displayed limit shows every one-sided derivative at $0$ exists and equals $0$, so all derivatives are continuous there.
\end{proof}

\begin{lemma}
\label{lem:2-21}
\lean{exists_one_zero_smooth_cutoff}
\uses{lem:2-20}
\leanok



\dcref{ch2:2.21}
For $r_1<r_2$, there is a smooth $h:\mathbb R\to\mathbb R$ with $h=1$ on $(-\infty,r_1]$, $0<h<1$ on $(r_1,r_2)$, and $h=0$ on $[r_2,\infty)$.
\end{lemma}

\begin{proof}
\leanok
\sketch With $f$ from Lemma~\ref{lem:2-20}, set
\[
h(t)=\frac{f(r_2-t)}{f(r_2-t)+f(t-r_1)}.
\]
The denominator is positive and the stated values follow from the support properties of $f$.
\end{proof}

\begin{definition}
\label{def:2-10-extra-1}
\lean{exists_one_zero_smooth_cutoff}
\uses{lem:2-21}
\leanok



A function with the properties of Lemma~\ref{lem:2-21} is a cutoff function.
\end{definition}

\begin{lemma}
\label{lem:2-22}
\lean{exists_smooth_ball_cutoff}
\uses{lem:2-21}
\leanok



\dcref{ch2:2.22}
For $0<r_1<r_2$, there is smooth $H:\mathbb R^n\to\mathbb R$ with $H=1$ on $\overline B_{r_1}(0)$, $0<H<1$ on $B_{r_2}(0)\setminus\overline B_{r_1}(0)$, and $H=0$ outside $B_{r_2}(0)$.
\end{lemma}

\begin{proof}
\leanok
\sketch Set $H(x)=h(|x|)$ for the cutoff $h$ of Lemma~\ref{lem:2-21}. It is smooth away from $0$ and constant near $0$.
\end{proof}

\begin{definition}
\label{def:2-10-extra-2}
\lean{SmoothBumpFunction}
\uses{lem:2-22}

\mathlibok

A smooth bump function is a smooth real-valued function equal to $1$ on a specified set and zero outside a specified neighborhood; Lemma~\ref{lem:2-22} supplies such functions on Euclidean balls.
\end{definition}

\begin{definition}
\label{def:2-10-extra-3}
\lean{HasCompactSupport.of_compactSpace}

\mathlibok
For a real- or vector-valued function $f$ on $M$, define
\[
\operatorname{supp}f=\overline{\{p\in M:f(p)\neq0\}}.
\]
The function is supported in $U$ when this support lies in $U$, and compactly supported when the support is compact. Every function on a compact space is compactly supported.
\end{definition}

\begin{definition}
\label{def:2-10-extra-4}
\lean{SmoothPartitionOfUnity.IsSubordinate}

\mathlibok
For an open cover $\mathcal X=(X_\alpha)_{\alpha\in A}$ of a topological space $M$, a partition of unity subordinate to $\mathcal X$ is a family of continuous functions $\psi_\alpha:M\to\mathbb R$ satisfying
\[
0\leq\psi_\alpha\leq1,\qquad \operatorname{supp}\psi_\alpha\subseteq X_\alpha,
\]
with locally finite supports and $\sum_{\alpha\in A}\psi_\alpha(x)=1$ for every $x$. Local finiteness makes the sum locally finite. It is smooth when every $\psi_\alpha$ is smooth.
\end{definition}

\begin{theorem}[Existence of Partitions of Unity]
\label{thm:2-23}
\lean{SmoothPartitionOfUnity.exists_isSubordinate}
\uses{lem:1-13, lem:2-22, prop:1-19}

\mathlibok

\dcref{ch2:2.23}
Every indexed open cover of a smooth manifold (with or without boundary) admits a smooth subordinate partition of unity.
\end{theorem}

\begin{proof}
\sketch For the boundaryless case, refine the cover by a countable locally
finite family of regular coordinate balls $B_i$ whose closures remain locally
finite. Choose enlarged coordinate balls $B_i'\subseteq X_{a(i)}$ and charts
$\varphi_i:B_i'\to B_{r_i'}(0)$ taking $\overline B_i$ to
$\overline B_{r_i}(0)$, where $r_i<r_i'$. Lemma~\ref{lem:2-22} gives $H_i$ positive on
$B_{r_i}(0)$ and zero outside it. Thus $H_i\circ\varphi_i$ on $B_i'$ extends
smoothly by zero outside $\overline B_i$ to a function $f_i$ supported on
$\overline B_i$. Then $f=\sum_i f_i$ is smooth and positive. The normalized
functions $g_i=f_i/f$ satisfy $\sum_i g_i=1$. Reindex by
$\psi_\alpha=\sum_{a(i)=\alpha}g_i$; local finiteness gives smoothness,
support in $X_\alpha$, and $\sum_\alpha\psi_\alpha=1$. The boundary case is
analogous using half-ball charts.
\end{proof}


\section{Applications of Partitions of Unity}

\begin{definition}
\label{def:2-11-extra-1}
\lean{Set.EqOn}

\mathlibok
If $A$ is closed in a topological space $M$ and $U$ is open with $A\subseteq U$, a bump function for $A$ supported in $U$ is a continuous $\psi:M\to\mathbb R$ with $0\leq\psi\leq1$, $\psi|_A\equiv1$, and $\operatorname{supp}\psi\subseteq U$.
\end{definition}

\begin{proposition}[Existence of Smooth Bump Functions]
\label{prop:2-25}
\lean{exists_smooth_bump_function_for}
\leanok



\dcref{ch2:2.25}
For every closed $A$ in a smooth manifold $M$ and every open $U\supseteq A$, there is a smooth bump function for $A$ supported in $U$.
\end{proposition}

\begin{proof}
\leanok
\sketch Apply Theorem~\ref{thm:2-23} to $\{U,M\setminus A\}$, obtaining $\psi_0,\psi_1$. Since $\psi_1$ vanishes on $A$, $\psi_0=1$ on $A$ and has support in $U$.
\end{proof}

\begin{proposition}[Existence of Smooth Exhaustion Functions]
\label{prop:2-28}
\lean{exists_positive_smooth_exhaustion_function}
\leanok


\dcref{ch2:2.28}
Every smooth manifold with or without boundary admits a positive smooth exhaustion function.
\end{proposition}

\begin{proof}
\leanok
\sketch Choose a countable cover by precompact open sets $V_j$ and a subordinate partition $\psi_j$. Set $f=\sum_{j\geq1}j\psi_j$. Local finiteness gives smoothness and $f\geq1$. For $N>c$, outside $\bigcup_{j=1}^N\overline V_j$ one has $f\geq N>c$; hence $f^{-1}(( -\infty,c])$ is closed in a finite union of compact closures and is compact.
\end{proof}

\begin{theorem}[Level Sets of Smooth Functions]
\label{thm:2-29}
\lean{exists_nonneg_smooth_zero_set_eq_of_isClosed}
\leanok


\dcref{ch2:2.29}
For every closed subset $K$ of a smooth manifold $M$, there is a smooth nonnegative $f:M\to\mathbb R$ with $f^{-1}(0)=K$.
\end{theorem}

\begin{proof}
\leanok
\sketch In $\mathbb R^n$, cover $\mathbb R^n\setminus K$ by countably many balls $B_{r_i}(x_i)$ with $r_i\leq1$. Let $h$ be a smooth bump equal to $1$ on $\overline B_{1/2}(0)$ and supported in $B_1(0)$. Choose $C_i\geq1$ bounding $h$ and its derivatives through order $i$, and define
\[
f(x)=\sum_{i=1}^\infty\frac{r_i^i}{2^iC_i}h\!\left(\frac{x-x_i}{r_i}\right).
\]
The series and each differentiated series converge uniformly by the $M$-test (the $i$th derivative terms are bounded by $2^{-i}$ for $i\geq k$), so $f$ is smooth, vanishes on $K$, and is positive off $K$. For general $M$, choose a coordinate-ball cover and subordinate partition $\psi_\alpha$; on each ball construct $f_\alpha\geq0$ with zero set $K\cap B_\alpha$, and use $f=\sum_\alpha\psi_\alpha f_\alpha$.
\end{proof}

\begin{exercise}
\label{exer:2-1}
\lean{ContMDiffMap.coe_mul}
\mathlibok
\dcref{ch2:2.1}
Let $M$ be a smooth manifold with or without boundary. Show that pointwise multiplication turns ${C}^{\infty }\left( M\right)$ into a commutative ring and a commutative and associative algebra over $\mathbb{R}$ . (See Appendix B, p. 624, for the definition of an algebra.)
\end{exercise}

\begin{exercise}
\label{exer:2-11}
\lean{contMDiff_subtype_val}
\mathlibok
\dcref{ch2:2.11}
\begin{itemize}
\item Exercise 2.11. Prove parts (a)-(c) of the preceding proposition.
\end{itemize}
\end{exercise}

\begin{exercise}
\label{exer:2-24}
\lean{SmoothPartitionOfUnity.exists_isSubordinate}
\mathlibok
\dcref{ch2:2.24}
Show how the preceding proof needs to be modified for the case in which $M$ has nonempty boundary.
\end{exercise}

\begin{exercise}
\label{exer:2-7}
\lean{ContMDiff.contMDiffOn}
\mathlibok
\dcref{ch2:2.7}
\begin{itemize}
\item Exercise 2.7. Prove the preceding two propositions.
\end{itemize}
\end{exercise}
